\documentclass[11pt]{article}

\usepackage[margin=1in]{geometry}
\usepackage{amsmath}
\usepackage{booktabs}
\usepackage{graphicx}
\usepackage[numbers]{natbib}
\usepackage[colorlinks=true,linkcolor=blue,citecolor=blue,urlcolor=blue]{hyperref}
\usepackage{microtype}

\newcommand{\pp}{\ensuremath{\,\mathrm{pp}}}
\title{When to Remember: A Preregistered Study of Failure-Memory\\Delivery Timing in an LLM Coding Agent}
\author{Joshua Warren\\Warren Applied Labs\\\texttt{joshua@warrenappliedlabs.com}}
\date{August 2026}

\begin{document}
\maketitle

\begin{abstract}
Long-term memory systems for LLM agents deliver recalled context at the start
of a turn, before the agent begins working. We ask whether that delivery point,
rather than memory content, limits the value of failure memories. In a
preregistered, hash-bound experiment on 30 synthetic TypeScript repair tasks
with deliberate trap fixes, one open-weight 35B coding model repeated its own
previously observed failure in 44.2\% of unassisted episodes. Injecting a
matched memory of that failure at turn start failed its registered support
test: repeated failures fell only 4.8 percentage points, and the task-pass
estimate favored an equivalent \emph{success} memory. Delivering the identical
fact at the moment the agent proposed the matched action eliminated
repetition: a confirmatory run on 18 preregistered tasks measured a
35.6-percentage-point reduction (95\% CI [16.7, 55.9], relative risk reduction
100\%, $p = 0.0019$), with zero invalid rows and zero protocol deviations.
The effect replicated across three measurements. Removing repeated failures
did not measurably improve task completion (+1.5\pp{} [0.0, 4.4]), so the
claim is failure-mode removal, not general improvement. Our first registration
returned NOT\_ESTIMABLE under its own zero-cut rule; we report both
registrations in full. The task corpus, harness, registrations, and analysis
artifacts are public.
\end{abstract}

\section{Introduction}
\label{sec:intro}

An agent with long-term memory will, sooner or later, face a situation it has
already gotten wrong once. What its memory system does in that moment decides
whether the stored experience has value. Long-term memory systems answer
the question the same way: retrieve whatever seems relevant and place it in
the prompt before the agent starts its turn
\citep{packer2023memgpt, chhikara2025mem0, xu2025amem, zhao2023expel}.
Retrieval quality, storage structure, and summarization policy vary widely
across these systems; the delivery point does not. Systems that do intervene
at tool-proposal time exist, but they deliver externally supplied safety
policies rather than the agent's own episodic memory
\citep{xiang2024guardagent, ruan2023toolemu}, and no prior work we know of
compares the two delivery points while holding the payload fixed.

This paper isolates the delivery point. We hold the memory itself fixed---the
same failure fact, the same wording, the same rendered token count---and vary
only \emph{when} it reaches the agent: at turn start, in the conventional
preamble position, or at the moment the agent proposes the specific action
that failed before. The distinction matters because coding agents fail in a
particular way: they take an attractive wrong path (silence the flaky-looking
test, edit the shadowed config file) that a general warning delivered minutes
of context earlier struggles to interrupt.

We built 30 synthetic TypeScript repair tasks around six classes of such
traps, gave a frozen open-weight model a first episode in which it fell into
the trap, then measured whether a memory of that failure prevented repetition
in a second, scored episode. The experiment was preregistered before data
collection: task splits, arms, seeds, outcome definitions, statistical
procedures, support thresholds, and decision rules were fixed in a
hash-bound registration that the harness verifies before any row executes.

Four results, in decreasing order of evidential strength:

\begin{enumerate}
\item \textbf{Action-site delivery eliminated repetition (confirmatory).} A
  matched failure fact delivered when the agent proposed the matched action
  reduced repeated failures by $35.6\pp$ (95\% CI $[16.7, 55.9]$, relative
  risk reduction 100\%, $p = 0.0019$) against turn-start delivery of the
  identical fact, on all 18 preregistered tasks with zero excluded data. The
  estimate and test operate on paired task-level means; descriptively, the
  pre-action arm repeated a known failure zero times in 270 valid episodes.
\item \textbf{Turn-start failure memory failed its registered test
  (confirmatory rejection).} Against a matched success memory it reduced
  repeated failures by only $4.8\pp$ $[0.7, 10.0]$ while its task-pass
  estimate was $6.3\pp$ \emph{worse} $[-15.6, +0.7]$; the registered
  compound test rejected it.
\item \textbf{Prevention did not become completion (bounded null).} The
  confirmatory run's task-pass benefit was $+1.5\pp$ $[0.0, 4.4]$
  ($p = 0.51$). The mechanism removes a known failure mode; it does not make
  the agent better overall.
\item \textbf{Preregistration bound us (methodological).} Our first
  registration ended NOT\_ESTIMABLE when a single cap-terminated episode
  produced an unclassifiable check and the registered rule allowed zero task
  cuts. The confirmatory result comes from a second registration whose
  scoring rule was fixed before any new data existed. We report both.
\end{enumerate}

The task corpus with its trap taxonomy, the execution harness, both
registration documents, the machine-readable decision rules, and the
per-episode logs and statistics artifacts of every registered run are public;
the analysis replays from the released episode logs byte-for-byte
(Section~\ref{sec:repro}). Per-attempt trace archives remain operator-local
with published hashes.

\section{Related work}
\label{sec:related}

\paragraph{Memory systems for LLM agents.}
Architectures such as MemGPT \citep{packer2023memgpt}, Mem0
\citep{chhikara2025mem0}, and A-MEM \citep{xu2025amem} manage what an agent
stores and retrieves across sessions; generative-agent simulacra
\citep{park2023generativeagents} and cognitive-architecture framings
\citep{sumers2023coala} organize retrieval around recency, importance, and
task structure. Benchmarks such as LoCoMo \citep{maharana2024locomo} measure
whether retrieved memories are \emph{correct}, and recent systems
characterization work taxonomizes these designs by construction and
retrieval cost \citep{omri2026agentmemory}. Situating our design on the
natural axes---memory source (the agent's own episode), trigger (fingerprint
match on a proposed action), delivery point (turn start vs.\ action site),
payload control (byte-matched fact), and outcome (repetition of a known
failure)---the delivery-point axis is the one this literature holds
constant: retrieved context lands in the prompt before generation. That axis
is our sole manipulation.

\paragraph{Learning from failure and experience.}
Reflexion stores verbal self-critiques of failed attempts for subsequent
trials \citep{shinn2023reflexion}; ExpeL distills cross-task insights from
experience pools \citep{zhao2023expel}; Voyager grows a skill library from
successes \citep{wang2023voyager}; Agent Workflow Memory induces reusable
workflows \citep{wang2024awm}; Self-Refine iterates on feedback within a task
\citep{madaan2023selfrefine}. These systems vary what is learned and stored.
Our contribution is orthogonal and narrower: given a single, already-correct
failure memory, we show its value depends on when it lands.

\paragraph{Position and salience effects in context.}
Language models use long contexts unevenly, privileging beginnings and ends
\citep{liu2024lost}, are measurably distracted by irrelevant context
\citep{shi2023distracted}, and are sensitive to seemingly neutral prompt
formatting \citep{sclar2023formatspread}. These results motivate our
hypothesis that a warning's distance from the decision it must influence is
load-bearing, but none of them tests intervention timing in an agentic loop
with a frozen, matched payload. We do not claim to isolate the mechanism;
Section~\ref{sec:discussion} treats candidate explanations as post hoc.

\paragraph{Coding agents and benchmarks.}
Agentic software engineering is standardly evaluated by end-to-end issue
resolution \citep{jimenez2023swebench, yang2024sweagent, wang2024openhands}
over ReAct-style tool loops \citep{yao2023react, schick2023toolformer}, with
code generation measured since HumanEval \citep{chen2021humaneval}. Repeating
one's own recorded failure is a distinct, narrower behavior that end-to-end
pass rates do not surface; our tasks are built to measure exactly it.

\paragraph{Preregistration and statistical rigor.}
Preregistration is established in medicine and psychology and has been argued
for in NLP \citep{vanmiltenburg2021preregistering} alongside broader
reproducibility programs \citep{pineau2021reproducibility} and calls for
uncertainty-aware evaluation \citep{miller2024errorbars}. Agent evaluations
in particular accumulate researcher degrees of freedom (caps, retries,
scoring judgments); our design commits every such choice to a hash-verified
document before execution, and our first registration's fate
(Section~\ref{sec:registrations}) shows the commitment binding.

\paragraph{Runtime guardrails at tool time.}
Sandboxed risk probing \citep{ruan2023toolemu} and guard agents that check
proposed actions against policies \citep{xiang2024guardagent} intervene at
the same point in the loop as our pre-action arm. That work targets safety
policies supplied from outside; we deliver the agent's own episodic failure
memory, advisory-only, and measure repetition of a specific known mistake.

\section{Experimental design}
\label{sec:method}

\subsection{Task corpus}
\label{sec:tasks}

We generated 30 self-contained TypeScript repair tasks, each a small
repository with a failing behavior, an executable functional contract, and a
deterministic offline checker that classifies any resulting repository state
as \texttt{PASS}, a specific failure class, or indeterminate. Each task
embeds one \emph{trap}: a wrong fix that is more attractive than the correct
one. Six trap classes, five tasks each:

\begin{table}[h]
\centering
\small
\begin{tabular}{ll}
\toprule
Trap class & The attractive wrong fix \\
\midrule
config-shadowing & edit the config file that is silently shadowed by another \\
flaky-looking-test & mark the failing check flaky instead of fixing the state bug \\
hidden-invariant & patch the symptom while violating an undocumented invariant \\
misleading-error-message & follow the error text to the wrong layer \\
stale-cache-illusion & ``fix'' code whose behavior is a stale cache artifact \\
wrong-layer-fix & patch the call site instead of the owning module \\
\bottomrule
\end{tabular}
\caption{Trap taxonomy. Within a class, no two tasks exceed 0.40 normalized
Jaccard similarity on state-defining sources or check scripts.}
\label{tab:traps}
\end{table}

Each task ships three run-time variants (surface renamings that preserve the
trap), a trapped revision, a correct-fix revision, and a no-trap control
revision. Splits were frozen before any registered run: 12 pilot tasks and 18
main tasks. There is no dev split; harness development used separately
generated tasks.

A registered trap audit gates every model profile: on all 30 tasks with no
memory assistance, the model must fall into the trap in at least 30\% of
episodes and fail to produce a correct fix in at least 50\%. The profile used
here audited at 40.0\% trapped both times it was audited (63.3\% and 70.0\%
non-fixed).

\subsection{Two-episode protocol and outcomes}
\label{sec:protocol}

Episode 1 lets the agent fail: it attempts the task, falls into the trap, and
the harness freezes that history---the failure fact text, a fingerprint of
the trapping action, and repository hashes. Episode 2 is the scored episode:
the agent faces the same task (per variant and seed) with an injected memory
that depends on the arm.

A scored episode counts as a \textbf{repeated failure} only when both hold:
(i) an episode-2 action matches the frozen trap fingerprint, and (ii) the
offline checker proves the same failure class as episode 1. Text similarity
alone never counts. \textbf{Task pass} requires the checker to return
\texttt{PASS} within all execution caps. Steps are counted from the action
trace.

\subsection{Arms and matched-injection contracts}
\label{sec:arms}

\begin{table}[h]
\centering
\small
\begin{tabular}{lll}
\toprule
Arm & Turn-start injection & Pre-action gate \\
\midrule
\textsc{no\_memory} & none & off \\
\textsc{turn\_start\_failure} & the frozen failure fact & off \\
\textsc{turn\_start\_success} & matched twin-repo success fact & off \\
\textsc{pre\_action\_failure} & none & advisory warning on matched action \\
\textsc{both} & the frozen failure fact & advisory warning on matched action \\
\bottomrule
\end{tabular}
\caption{Arms. \textsc{both} is registered as secondary and enters no primary
test.}
\label{tab:arms}
\end{table}

The \textbf{timing pair} (\textsc{turn\_start\_failure} vs.\
\textsc{pre\_action\_failure}) delivers the same single fact with identical
fact ID, citation hash, and rendered token count; only the frame differs as
needed to place it at the two delivery points. The gate watches proposed
actions; when a proposal matches the frozen fingerprint, it injects the fact
as an advisory and lets the model choose again. On any gate error the action
runs without a note, and the registered rules invalidate such an episode
(this never occurred). Delivery at the action site is therefore a
\emph{package}: the matched proposal is intercepted rather than executed,
the fact arrives as a direct response to it, and the model re-decides---one
additional model interaction on each match. That package is inseparable from
the delivery point (a warning about an action cannot arrive ``at'' the
action after it has executed), and the registered hypothesis compares
delivery packages, not a pure clock shift. The design does not include a
neutral-interruption control (an interception carrying no memory), so it
does not isolate the marginal contribution of the fact's content within the
package (Section~\ref{sec:threats}).

The \textbf{content pair} (\textsc{turn\_start\_failure} vs.\
\textsc{turn\_start\_success}) matches its two facts on path-shape and
action-shape hashes, token-set Jaccard in $[0.80, 1.00]$, and a rendered
token gap of at most 8 tokens and 5\%.

Every arm of a cell starts from the same task commit, memory image, tools,
caps, seed, and model profile, with per-arm isolated worktrees, memory
directories, and sessions; start-state hashes are recorded per row and
audited for equality within cells.

\subsection{Model and execution}
\label{sec:model}

All episodes run one frozen open-weight model profile: the Qwen-family
\texttt{qwen3.5:35b} release (Q4\_K\_M quantization, 32{,}768 token
context, temperature 0, no thinking mode, 4{,}096 max output tokens). The
model is identified as an artifact by its served-model digest
(\texttt{3460ffee\ldots}) and profile hash, re-verified before every run;
no technical report for this release had been published at the time of
writing, so we cite the closest published family report
\citep{qwen2024qwen25} for architectural context only. Caps per episode: 12 turns, 8 tool calls,
20{,}480 cumulative tokens, 600\,s duration, 180\,s per request. Host and API
faults retry up to five times after the first try; exhaustion pauses the run
rather than voiding a row; a returned task result, valid or invalid, is never
re-executed.

\subsection{Statistical analysis}
\label{sec:stats}

The task is the analysis unit: row outcomes are averaged over variants and
seeds within each task and arm, and every test operates on the 18 paired
task means. Intervals are 10{,}000-draw grouped percentile bootstraps (95\%
for primaries, 90\% for the no-trap diagnostic); $p$-values come from
10{,}000-draw paired sign-flip randomization tests; the analysis seed (81)
is registered. The content hypothesis uses a compound
$p = \max(p_{\text{repeated-failure}}, p_{\text{task-pass}})$. Holm
correction spans the registered family: two members in the first
registration, one in the second. The timing support rule requires an
absolute repeated-failure benefit of at least $5\pp$, a point relative risk
reduction of at least 30\%, an interval lower bound strictly above zero, and
adjusted $p < 0.05$. A pilot on the disjoint 12-task split simulated 0.8364
power for the timing test at 18 tasks (0.80 required). The pilot simulation
used the first registration's two-member Holm family; the second
registration's single-member family can only make the same test easier, so
the transferred power estimate is conservative.

\subsection{Two registrations}
\label{sec:registrations}

\textbf{Registration 1} (decision rule v12) ran the full five-arm design plus
no-trap controls: $18 \times 3 \times 5 \times 7 = 1{,}890$ rows on seeds
1--5. It allowed zero primary task cuts: any invalid row in a compared cell
cuts the whole task and voids the affected hypothesis. All 1,890 rows
completed; 1,888 were valid. One of the two invalid rows---a pre-action
episode that hit the cumulative token cap mid-edit, leaving a repository
state the checker could not classify---sat in the timing comparison. The rule
fired as written: H6-timing was NOT\_ESTIMABLE despite a $+38.0\pp$
$[18.0, 59.6]$ exploratory estimate on the 17 complete tasks. The content
hypothesis, unaffected, was decided (rejected). The registered retry rule
forbids re-running any returned result, so a confirmatory timing answer
required new data.

\textbf{Registration 2} (decision rule v13) was fixed before any new row
executed. It decides H6-timing alone on the same 18 tasks with the two timing
arms and \emph{new} seeds 6--10 (540 rows). Three changes, all registered in
advance: a cap-driven unclassifiable check is scored worst-case against the
pre-action arm (repeated failure in \textsc{pre\_action\_failure}, no
repeated failure and a pass in \textsc{turn\_start\_failure}) instead of
cutting the task; the hypothesis family has one member; and the pilot's power
evidence transfers by pinned artifact hashes, with the harness required to
reproduce the pilot's statistics byte-for-byte before any row runs. The
zero-cut rule stands for every other invalid reason. The worst-case scoring
rule can only shrink the measured effect, and in the event it was never
exercised: the confirmatory run had zero invalid rows.

Registration 2 executed twice, and we report both executions. The first
manifest completed all 540 episodes and then crashed in the harness's
post-run statistics step (an option-plumbing defect caught by the analysis
guard). Because runs are hash-bound to the harness that started them, that
run directory was abandoned with a deviation record, and the registration
was re-executed on the repaired, test-covered harness with the same
registered seed set. The abandoned manifest's episode log had already been
written by the crashed finalization attempts, so raw arm-level outcomes
existed on disk; they were never statistically analyzed before the
re-execution decision, but during supervision we observed aggregate
final-state counts without arm attribution, and a skeptic need not take
non-observation on faith. We therefore publish both episode logs and report
both descriptively: the abandoned manifest shows 100/270 turn-start repeated
failures against 0/270 pre-action ($+37.0\pp$ episode-level); the kept,
confirmatory manifest shows 96/270 against 0/270 ($+35.6\pp$). The kept
manifest's effect is the \emph{smaller} of the two, so selection between
manifests cannot manufacture the result, and the conclusion is identical
under either log. The registered analysis ran exactly once, on the final
manifest.

\section{Results}
\label{sec:results}

\subsection{Primary outcomes}
\label{sec:primary}

\begin{table}[t]
\centering
\small
\begin{tabular}{llrrrrl}
\toprule
Comparison & Status & Tasks & RF benefit & Task-pass benefit & $p$ (adj.) & Decision \\
\midrule
Timing (R2) & confirmatory & 18/18 & $+35.6\pp$ $[16.7, 55.9]$ & $+1.5\pp$ $[0.0, 4.4]$ & 0.0019 & \textbf{supported} \\
Content (R1) & confirmatory & 18/18 & $+4.8\pp$ $[0.7, 10.0]$ & $-6.3\pp$ $[-15.6, +0.7]$ & 0.9393 & rejected \\
Timing (R1) & exploratory & 17/18 & $+38.0\pp$ $[18.0, 59.6]$ & --- & --- & not estimable \\
\bottomrule
\end{tabular}
\caption{Primary results. RF benefit is the reduction in task-level repeated
failure rate for the candidate arm (pre-action delivery in timing rows;
turn-start failure vs.\ turn-start success in the content row), with 95\%
bootstrap intervals. The content $p$ is the registered compound value; the
R1 timing row is exploratory by the registered zero-cut rule; because its
decision is not estimable, no $p$-value is shown for it (the unadjusted
descriptive sign-flip statistic on the 17 complete tasks was 0.0018).}
\label{tab:primary}
\end{table}

Table~\ref{tab:primary} contains the study. All intervals and $p$-values
operate on the 18 paired task-level means; episode counts are descriptive.
In the confirmatory run the pre-action arm repeated a known failure
\textbf{zero} times in 270 valid episodes, against a 35.6\% task-level
repetition rate for turn-start delivery of the identical fact---a relative
risk reduction of 100\% (95\% CI $[100, 100]$, degenerate because the
candidate numerator is zero in every bootstrap draw). Every registered support
condition held with margin. The estimate is consistent across three
measurements: the pilot ($+37.8\pp$, descriptive), the first registration's
exploratory estimate ($+38.0\pp$ on 17 complete tasks), and the confirmatory
run ($+35.6\pp$ on all 18).

The content comparison is a genuine negative with structure. Turn-start
failure wording did reduce repetition relative to matched success wording
($+4.8\pp$, interval above zero), but the registered rule demanded joint
support on completion as well, and there the failure fact \emph{cost}
$6.3\pp$ against the success fact. Under the compound test the hypothesis
was rejected. Placement, not content, is what our data supports changing.

\subsection{Descriptive arm outcomes}
\label{sec:descriptive}

\begin{table}[t]
\centering
\small
\begin{tabular}{lrrrr}
\toprule
Arm & Valid rows & Repeated failures & RF rate & Pass rate \\
\midrule
\multicolumn{5}{l}{\emph{Registration 1 (seeds 1--5)}} \\
\textsc{no\_memory} & 269 & 119 & 44.24\% & 2.60\% \\
\textsc{turn\_start\_success} & 270 & 110 & 40.74\% & 7.04\% \\
\textsc{turn\_start\_failure} & 270 & 97 & 35.93\% & 0.74\% \\
\textsc{pre\_action\_failure} & 269 & 0 & 0.00\% & 2.97\% \\
\textsc{both} & 270 & 0 & 0.00\% & 1.48\% \\
\midrule
\multicolumn{5}{l}{\emph{Registration 2 (seeds 6--10)}} \\
\textsc{turn\_start\_failure} & 270 & 96 & 35.56\% & 1.85\% \\
\textsc{pre\_action\_failure} & 270 & 0 & 0.00\% & 3.33\% \\
\bottomrule
\end{tabular}
\caption{Raw episode-level outcomes before task-level aggregation,
descriptive only. Pass rates are low in every arm because the registered
caps bind hard (Section~\ref{sec:threats}).}
\label{tab:arms-outcomes}
\end{table}

Table~\ref{tab:arms-outcomes} shows the unaggregated picture. Three
observations. First, the unassisted repetition rate (44.2\%) confirms the
traps do what they were built to do. Second, pre-action delivery drove
repetition to zero in both registrations, on disjoint seeds. Third, the
highest raw pass rate belongs to turn-start \emph{success} memory (7.04\%),
and \textsc{both}---gate plus turn-start failure fact---passed less
(1.48\%) than the gate alone (2.97\%): whatever failure wording does to an
agent's preamble, it is not obviously free even when the gate handles
prevention.

\subsection{Prevention is not completion}
\label{sec:completion}

The confirmatory task-pass benefit of pre-action delivery is $+1.5\pp$ with
an interval touching zero ($[0.0, 4.4]$, $p = 0.51$). Steps were essentially
unchanged (5.80 vs.\ 5.70 mean steps per episode). Blocking the remembered
wrong path does not, in this corpus and under these caps, convert into
solving the task: agents typically moved from the trap to a different
unsuccessful approach rather than to the correct fix.

\subsection{No-trap diagnostic}
\label{sec:notrap}

Registration 1 also measured whether the gate changes behavior when no
matching trap exists (540 no-trap control rows). The steps difference was
$-0.026$ with a 90\% interval $[-0.063, +0.0037]$, well inside the
registered $\pm 2$ margin: no detectable hesitation. The pass-rate half of
the diagnostic is recorded as \emph{unresolved}: its 90\% interval
$[-0.74, +2.22]\pp$ misses the registered $\pm 2\pp$ margin by $0.22\pp$,
and a registered amendment documents that this margin is finer than the
design can resolve at any feasible task count. The point estimate favors the
gate arm ($+0.74\pp$); nothing in the data suggests harm.

\subsection{Cost}
\label{sec:cost}

The two registered runs consumed 56.1M tokens of episode traffic (43.4M and
12.7M). All statistics are computed deterministically from row artifacts
with zero model calls, and replay byte-identically from the immutable
episode logs.

\section{Threats to validity}
\label{sec:threats}

\paragraph{Construct.}
Repeated failure is defined conservatively (fingerprint match \emph{and}
checker-proven same failure class), so our rates understate softer forms of
repetition. The traps are synthetic; they compress into one file tree what
real repositories spread across services. The two-episode protocol hands the
gate a correct, matched memory; production recall is noisier, and our effect
is conditional on the match being right.

\paragraph{Internal.}
The timing pair pins fact identity, citation hash, and rendered token count;
arms are isolated per worktree, memory, and session, with audited start-state
equality. The gate is advisory and fail-open, so the treatment cannot operate
by blocking. The abandoned first manifest of registration 2 is handled by evidence
rather than assertion: both episode logs are published, the registered
analysis ran only on the final manifest, and the abandoned manifest's
descriptive effect is larger than the kept one's
(Section~\ref{sec:registrations})---so the only selection story available
runs against the reported result. The delivery-package structure of the
pre-action arm (interception plus re-decision, Section~\ref{sec:arms}) is a
construct limitation: a neutral-interruption control was not run, so the
package's effect is established but the marginal effect of the memory
content within it is not.

\paragraph{Statistical.}
Tasks are the unit, respecting within-task clustering across variants and
seeds; tests are exact permutation analogues; the confirmatory family has one
member and its $p$-value survives any correction. The zero-numerator relative
risk reduction makes the RRR interval degenerate at 100\%; the absolute
benefit interval $[16.7, 55.9]\pp$ is the informative bound.

\paragraph{External.}
One model, one quantization, one language, one task style, frozen caps. The
binding token cap suppresses pass rates in every arm and limits what
completion effects could be observed. Nothing here licenses claims about
other models, real repositories, or memory systems whose retrieval quality
differs from our matched setting. The corpus is public as of August 2026 and
carries a contamination canary; results on models trained after that date
require a contamination check before comparison.

\section{Discussion}
\label{sec:discussion}

\paragraph{Why placement might dominate.}
Three post hoc candidates, none isolated by our design. Distance: long-context
studies show mid-context material is used unevenly \citep{liu2024lost}, and a
turn-start warning is maximally distant from the decision it must influence.
Competition: by proposal time the context is dominated by freshly read code
and tool output, and irrelevant-context effects \citep{shi2023distracted}
suggest the warning loses that competition. Salience: the gate delivers the
fact as a direct response to the specific proposed action, when the
alternative is concrete rather than hypothetical. Separating these requires
designs that vary injection distance continuously; our corpus supports such
follow-ups.

\paragraph{Design implication for memory systems.}
Failure memories should be routed to a just-in-time advisory delivered
against the proposed action---the position where guard agents already sit
\citep{xiang2024guardagent}---rather than into the retrieval preamble.
Semantic and preference context can stay at turn start; nothing here argues
against preamble recall in general. The implication is a routing rule, not a
new architecture: store failure memories with an action fingerprint, match at
tool-proposal time, inject advisory-only, fail open. Our harness's gate is a
reference implementation. The rule is established only in the tested regime:
a correct, byte-matched memory of this task's own failure, matched by an
exact fingerprint, on synthetic tasks. Production systems face false and
missed matches, noisy retrieval, advisory fatigue under repeated triggers,
and per-proposal matching cost; none of these is measured here, and each
could erase the benefit.

\paragraph{What preregistration bought.}
The first registration's NOT\_ESTIMABLE outcome is the strongest evidence we
can offer that the protocol binds: a $+38\pp$ effect with $p = 0.0036$ was
set aside because one row violated a rule we wrote ourselves, and the fix was
registered before new data existed rather than argued after the fact. The
cost was one additional run. For agent evaluations---where caps, retries, and
scoring judgments hide many researcher degrees of freedom---we found the
discipline cheap relative to what it protects, echoing arguments made for
NLP broadly \citep{vanmiltenburg2021preregistering, pineau2021reproducibility}.

\section{Reproducibility and data availability}
\label{sec:repro}

The task corpus (30 tasks, trap taxonomy, JSON schemas, frozen inventory hash
\texttt{687615b5\ldots}), the execution harness, both preregistration
documents, the machine-readable decision rules (v12, v13), and the study
report are public in the project repository at tag
\texttt{h6-study-2026-08}.\footnote{\url{https://github.com/joshuaswarren/remnic/tree/h6-study-2026-08}
--- harness under \texttt{packages/bench/}, registrations and report under
\texttt{docs/research/failure-gate/}, corpus under
\texttt{packages/bench/fixtures/h6-failure-gate/}.} The per-episode logs,
run manifests, decision-rule copies, deviation logs, and statistics, power,
and audit artifacts of all four runs (both registrations, the bound pilot,
and the abandoned manifest) are published as release assets on the same
tag.\footnote{\url{https://github.com/joshuaswarren/remnic/releases/tag/h6-study-2026-08}}
The registered analysis replays from the released episode logs byte-for-byte
with zero model calls via a single command. Per-attempt trace archives (full
agent transcripts) remain operator-local; their SHA-256 hashes are bound in
the released manifests. One disclosure: the confirmatory run's replay executes under
a harness newer than the one that ran the episodes (post-run report-layer
fixes); the replay verifies the decision artifacts byte-for-byte and records
both harness hashes side by side.

\section*{Acknowledgments}
The experiment harness, execution supervision, and analysis tooling were
built with substantial assistance from LLM coding agents operating under the
author's direction, and every number in this paper was verified by
deterministic replay against the released artifacts rather than taken from
agent output. The study design, registrations, and claims are the author's.
No external funding.

\bibliographystyle{plainnat}
\bibliography{references}

\appendix

\section{Amendment log}
\label{app:amendments}

Registration 1 carried three amendments, each registered before the data it
governed: (2) the cumulative token cap rose from 16{,}384 to 20{,}480 after
the pilot showed the original cap consumed by re-sent context before any
task could complete, with the replacement value selected by four 30-task
trap audits rather than by outcome data; (3) the content hypothesis was
removed from the launch power gate after two pilots measured its power at
0.000 (it remained registered, estimated, and decided); (4) the no-trap
pass-rate margin was removed from the launch gate after power analysis
showed no feasible task count could resolve it (the estimate remains
reported, Section~\ref{sec:notrap}). Registration 2's changes are listed in
Section~\ref{sec:registrations}. Full texts ship in the repository.

\section{Integrity receipts}
\label{app:receipts}

Every registered run writes a frozen manifest (expected row keys and order),
per-try checkpoints, an append-only deviation log, and hash-bound statistics,
power, and audit artifacts; the paper report regenerates only from these.
Registration 1: 1,890 rows, deviation log empty, statistics
\texttt{16fd5af1\ldots}. Registration 2 confirmatory run
(\texttt{h6-51d25af1ed731c35a94c28fe}): 540 rows, deviation log empty,
statistics \texttt{0cedc1b8\ldots}, power \texttt{2f1630c7\ldots}, audit
\texttt{38e14e5a\ldots}. Pilot transfer: manifest \texttt{5bce2c0c\ldots},
power \texttt{cfaec22f\ldots}. Full 64-character hashes are in the
repository's study report.

\end{document}
